% Generated draft: requires mathematical and source review.
\documentclass[11pt]{article}
\usepackage[margin=25mm]{geometry}
\usepackage{fontspec}
\setmainfont{Latin Modern Roman}
\usepackage{amsmath,amssymb,mathtools,unicode-math}
\setmathfont{Latin Modern Math}
\usepackage{enumitem,needspace,xurl,hyperref}
\setlength{\parindent}{0pt}
\setlength{\parskip}{5pt}
\setlength{\emergencystretch}{3em}
\newcommand{\N}{\mathbb{N}}
\newcommand{\Z}{\mathbb{Z}}
\newcommand{\Q}{\mathbb{Q}}
\newcommand{\R}{\mathbb{R}}
\newcommand{\C}{\mathbb{C}}
\newcommand{\F}{\mathbb{F}}
\newcommand{\A}{\mathbb{A}}
\newcommand{\PP}{\mathbb{P}}
\newcommand{\E}{\mathbb{E}}
\newcommand{\eps}{\varepsilon}
\newcommand{\dd}{\,\mathrm{d}}
\begin{document}
\section*{UnsolvedMath ID: 10600040}
\textbf{Source problem code:} catalogue entry
\subsection{Definitions and mathematical statement}
Let \(D\) be a finite chord diagram on a circle. Two chords are linked when their endpoints alternate around the circle. Its intersection graph \(G_D\) has chords as vertices and linked pairs as edges. Define the independence complex
\[
 \mathcal I(D)=\{S\subseteq V(G_D):S\text{ consists of pairwise unlinked chords}\}.
\]
A \(k\)-element set is a simplex of dimension \(k-1\), including the empty simplex in dimension \(-1\). Fix an order on the chords. The source's degree-raising differential can be written as the augmented simplicial \emph{cochain} differential
\[
 \delta[S]=\sum_{a\notin S\, :\,S\cup\{a\}\in\mathcal I(D)}
 (-1)^{|\{s\in S:s<a\}|}[S\cup\{a\}],
\]
with \(\delta^2=0\) over \(\Z\). The source calls this a chain complex although its differential adds a chord; dualizing gives the usual boundary convention. When \(G_D\) consists of \(n\) disjoint edges, \(n\geq1\), the complex is \((S^0)^{*n}\cong S^{n-1}\), as in the supplied example.

The first question asks whether chord diagrams can produce complexes other than bouquets of spheres. The homotopy-theoretic reading is whether some \(|\mathcal I(D)|\) is not homotopy equivalent to a wedge of spheres; the source does not explicitly distinguish this from a homeomorphism or combinatorial-isomorphism question, or fix disconnected/empty-complex conventions.

The second question concerns diagrams of classical or virtual knots \emph{and links} whose all-\(A\) smoothing has exactly one circle. For such a diagram, the crossings give a chord diagram on that state circle. Find combinations of Reidemeister moves such that whenever two such diagrams represent equivalent knots or links, they can be joined by a chain of the allowed moves with every intermediate diagram still having a single-circle all-\(A\) state, and with the corresponding chord-diagram homology and homotopy unchanged along the chain. The elementary versus composite move convention, the virtual equivalence, and any grading shifts or weaker comparison notion must be specified.

The source motivates these questions by the reported identification of lowest-quantum-grading Khovanov homology with the above complex for single-circle all-\(A\) diagrams, and by Khovanov spectrification. It also discusses a flag-type triangulation condition: every clique in the one-skeleton spans a simplex. If that one-skeleton is realizable as a chord non-intersection graph, its clique complex is \(\mathcal I(D)\). Not every graph is a chord intersection or non-intersection graph, which motivates passing to graph-link theory; the source reports that general manifold examples then appear in lower-degree homology and suggests lower homotopy invariants. These are contextual reports, not claims that every manifold is represented by a chord diagram. The precise Khovanov/graph-link models and grading conventions are absent.
\subsection{Short English statement}
Ask which homotopy types occur as independence complexes of chord intersection graphs, beyond wedges of spheres, and find moves connecting equivalent single-circle all-\(A\)-state diagrams while preserving their associated homology and homotopy.
\subsection{Sources}
\begin{itemize}
\item \emph{Unsolved Problems in Virtual Knot Theory and Combinatorial Knot Theory}, revised survey, problem-list item 40. \url{https://arxiv.org/abs/1409.2823}. Recorded access 2026-07-29; extraction source-tex.
\item The statement refers to work with Lipshitz, later spelled Lipschitz, on Khovanov spectrification and to a Bloom--Nikonov approach, without complete references. The supplied research summary mentions Lipshitz--Sarkar spectrification and graph-link theory, but no dedicated resolving paper.
\item Local UnsolvedMath registry, \texttt{data/unsolvedmath/problems.json}, permanent ID \(10600040\), exact source code \texttt{catalogue entry}; supplied OPEN-TRIAGE status.
\end{itemize}
\end{document}
